\begin{table}[!h]
\centering
\caption{\label{tab:treat_status}Hypothetical Quota Assignment Across Election Cycles}
\centering
\fontsize{8}{10}\selectfont
\begin{threeparttable}
\begin{tabular}[t]{ccccc}
\toprule
GP & 2005 & 2010 & 2015 & 2020\\
\midrule
1 & W & O & W & O\\
2 & O & W & O & W\\
3 & W & W & O & O\\
4 & O & O & W & W\\
5 & W & O & O & W\\
\addlinespace
6 & O & W & W & O\\
7 & W & W & W & O\\
8 & O & O & O & W\\
9 & W & O & W & W\\
10 & O & W & O & O\\
\addlinespace
11 & W & W & O & W\\
12 & O & O & W & O\\
13 & W & O & O & O\\
14 & O & W & W & W\\
15 & W & W & W & W\\
\addlinespace
16 & O & O & O & O\\
\bottomrule
\end{tabular}
\begin{tablenotes}
\item \textit{Note: } 
\item W = Quota seat (reserved for women); O = Open seat (anyone can contest). This table illustrates the randomized assignment of gender quotas across four electoral cycles for a hypothetical set of 16 Gram Panchayats.
\end{tablenotes}
\end{threeparttable}
\end{table}
